\documentclass{article}
\begin{document}

\title{IDAxiBrillBH}
\author{Paul Walker, Steve Brandt}
\date{1 September 1999}
\maketitle

\abstract{Thorn IDAxiBrillBH provides analytic initial data for a vacuum
  black hole spacetime:  a single Schwarzschild black hole in
  isotropic coordinates plus Brill wave.  This initial data is
  provided for the 3-conformal metric, it's spatial derivatives, and
  extrinsic curvature.}

% The above author & date lines seem to be ignored, so for now I'll try to work
% around this

\section{Authors \& Date}
Paul Walker and Steve Brandt\\
1 September 1999

\section{Purpose}

The pioneer, Bernstein, studied a single black hole which is
non-rotating and distorted in azimuthal line symmetry of 2 dimensional 
case \cite{Bernstein93a}. In this non-rotating case, one chooses the
condition, $K_{ij} = 0$,  and 
\begin{equation}
\gamma_{ab} = \psi^4 \hat \gamma_{ab},
\end{equation}
where $\gamma_{ab}$ is the physical three metric and
$\hat{\gamma}_{ab}$ is some chosen conformal three metric.

The Hamiltonian constraint reduces to
\begin{equation}
\hat \Delta \psi = \frac{1}{8}\psi \hat R,
\label{eqn:conformal_hamiltonian}
\end{equation}
where $\hat \Delta$ is the covariant Laplacian and $\hat R$ is Ricci
tensor for the conformal three metric. This form allows
us to choose an arbitrary conformal three metric, and then solve an
elliptic equation for the conformal factor, therefore satisfying the
constraint equations ($K_{ij} = 0$ trivially satisfies the momentum
constraints in vacuum).  This approach was used to create
``Brill waves'' in a spacetime without black holes \cite{Brill59}.
Bernstein extended this to the black hole spacetime.  Using
spherical-polar coordinates, one can write the 3-metric,
\begin{equation}
\label{eqn:sph-cood}
ds^2 = \psi^4 (e^{2q} (dr^2 + r^2 d \theta^2) + r^2 \sin \theta d
\phi^2),
\end{equation}
where $q$ is the Brill ``packet'' which takes some functional form.
Using this ansatz with (\ref{eqn:conformal_hamiltonian}) leads to
an elliptic equation for $\psi$ which must be solved
numerically. Applying the isometry condition on $\psi$ at a finite
radius, and applying $M/2r$ falloff conditions on $\psi$ at the
outer boundary (the ``Robin'' condition), along with a packet which
obeys the appropriate symmetries (including being invariant under the
isometry operator), will make this solution describe a black hole with
an incident gravitational wave.  The choice of $q=0$ produces the
Schwarzschild solution.   The typical $q$ function used in
axisymmetry, and considered here in the non-rotating case, is 
\begin{equation}
q = Q_0 \sin \theta^n \left [ \exp\left(\frac{\eta -
      \eta_0^2}{\sigma^2}\right ) + \exp\left(\frac{\eta +
      \eta_0^2}{\sigma^2}\right ) \right ].
\end{equation}
Note regularity along the axis requires that the exponent $n$ must be
even.  Choose a logarithmic radial coordinate $\eta$, which is
related to the
asymptoticlly flat coordinate $r$ by $\eta = ln (2r/m)$, where $m$ is a
scale parameter.  One can rewrite (\ref{eqn:sph-cood}) as
\begin{equation}
ds^2 = \psi(\eta)^4 [ e^{2 q} (d \eta^2 +  d\theta^2) +  \sin^2
\theta d\phi^2].
\end{equation}

In the previous Breinstein work, the above $r$ is transformed to a
logarithmic radial coordinate

\begin{equation}
\label{eta_coord}
\eta = \ln{\frac{2r}{m}}.
\end{equation}

The scale parameter $m$ is equal to the mass of the Schwarzschild
black hole, if $q=0$.  In this coordinate, the 3-metric is
\begin{equation}
\label{eqn:metric_brill_eta}
ds^2 = \tilde{\psi}^4 (e^{2q} (d\eta^2+d\theta^2)+\sin^2 \theta
d\phi^2),
\end{equation}
and the Schwarzschild solution is 
\begin{equation}
\label{eqn:psi}
\tilde{\psi} = \sqrt{2M} \cosh (\frac{\eta}{2}).
\end{equation}
We also change the notation of $\psi$ for the conformal factor is same 
as $\tilde{\psi}$ \cite{Camarda97a}, for the $\eta$ coordinate has the
factor $r^{1/2}$ in the conformal factor.  Clearly $\psi(\eta)$ and
$\psi$ differ by a factor of $\sqrt{r}$.  The Hamiltonian 
constraint is
\begin{equation}
\label{eqn:ham}
\frac{\partial^2 \tilde{\psi}}{\partial \eta^2} + \frac{\partial^2
  \tilde{\psi}}{\partial \theta^2} + \cot \theta \frac{\partial
  \tilde{\psi}}{\partial \theta} = - \frac{1}{4} \tilde{\psi}
(\frac{\partial^2 q}{\partial \eta^2} + \frac{\partial^2 q}{\partial
  \theta^2} -1).
\end{equation}

For solving this Hamiltonian constraint numerically.  At first
we substitute
\begin{eqnarray}
\delta \tilde{\psi} & = & \tilde{\psi}+\tilde{\psi}_0 \\
                    & = & \tilde{\psi}-\sqrt{2m} \cosh(\frac{\eta}{2}).
\end{eqnarray}
to the equation~(\ref{eqn:ham}), then we can linearize it as
\begin{equation}
\frac{\partial^2 \delta\tilde{\psi}}{\partial \eta^2} + \frac{\partial^2
  \delta\tilde{\psi}}{\partial \theta^2} + \cot \theta \frac{\partial
  \delta\tilde{\psi}}{\partial \theta} = - \frac{1}{4}
(\delta\tilde{\psi} + \tilde{\psi}_0) (\frac{\partial^2 q}{\partial
  \eta^2} + \frac{\partial^2 q}{\partial \theta^2} -1).
\label{eqn:ham_linear}
\end{equation}
For the boundary conditions, we use for the inner boundary condition
an isometry condition:
\begin{equation}
\frac{\partial \tilde{\psi}}{\partial \eta}|_{\eta = 0} = 0,
\end{equation}
and outer boundary condition, a Robin condition:
\begin{equation}
(\frac{\partial \tilde{\psi}}{\partial \eta} + \frac{1}{2}
\tilde{\psi})|_{\eta=\eta_{max}} = 0.
\end{equation}

%  [[ DPR: What is this: ?? ]]
%This thorn provides
% \begin{enumerate}
%  \item CactusEinstein
% \end{enumerate}

\section{Comments}

We calculate equation~(\ref{eqn:ham_linear}) with spherical
coordinates.  However, Cactus needs Cartesian coordinates.  Therefore,
we interpolate $\psi$ to the Cartesian grid by using an interpolator.
Note that the interpolator has linear, quadratic, and cubic
interpolation.

% Automatically created from the ccl files:
\include{interface}
\include{param}
\include{schedule}

\bibliographystyle{prsty}
\begin{thebibliography}{10}
\bibitem{Bernstein93a}
  D. Bernstein, Ph.D thesis University of Illinois Urbana-Champaign,
  (1993)
\bibitem{Brill59}
  D. S. Brill,Ann. Phys.{\bf 7}, 466 (1959)
\bibitem{Camarda97a}
  K. Camarda, Ph.D thesis University of Illinois Urbana-Champaign, (1998)
\end{thebibliography}
\end{document}
